Welcome to the 11th Social Media Mining for Health Research and Applications and the Health Real-World Data (\#SMM4H-HeaRD) Workshop and Shared Tasks, co-located with the 64th Annual Meeting of the Association for Computational Linguistics (ACL 2026). This year’s workshop is held online, connecting the global community of researchers who develop and evaluate natural language processing, machine learning, and artificial intelligence methods for health-related text.
\par\vspace{1em}

Since its inception in 2016, \#SMM4H has grown from a small, focused workshop, into one of the field’s most visible recurring venues for health NLP research. Over eleven editions, the workshop has brought together hundreds of teams from across the world, spanning academia, industry, government, and clinical institutions. Participation has grown substantially year over year. Collectively, the \#SMM4H shared tasks and associated publications have accumulated over 1,700 citations, reflecting the community’s sustained reliance on the datasets, baselines, and benchmarking frameworks the workshop has produced.
\par\vspace{1em}

The \#SMM4H-HeaRD 2026 edition marks a record in the scope of our shared task program, with eight shared tasks spanning an unprecedentedly broad range of domains and data modalities: detection of adverse drug events in multilingual social media posts, detection of insomnia in clinical notes, estimating flu vaccine effectiveness from social media, generation of structured medical notes from dialogues, detection of patient metadata in SARS-CoV-2 sequencing articles, predicting TNM staging from pathology reports, extraction of social and clinical impacts of substance use, and multilingual clinical entity annotation projection and extraction. This year, 110 registered teams from 31 countries registered, and 79 teams participated in at least one task — the largest participation in the workshop’s history. The shared tasks included 8 tasks: 12 teams for Task 1, 8 for Task 2, 6 for Task 3, 5 for Task 4, 10 for Task 5, 7 for Task 6, 10 for Task 7, and 21 for Task 8. Winning systems were invited for oral presentations, and other accepted system description papers were included as poster presentations.In all, these proceedings include 53 papers.

\par\vspace{1em}
This year’s edition also highlights the field’s evolving methodological landscape. Large language models featured prominently across submissions, used for classification, information extraction, data augmentation, and structured output generation — reflecting both their growing accessibility and the community’s interest in evaluating their reliability in high-stakes health contexts. In alignment with the ACL 2026 Theme Track on the Explainability of NLP Models, we specifically invited work that advances the transparency, interpretability, and trustworthiness of models applied in clinical and public health settings.

\par\vspace{1em}
This year’s workshop features two invited keynotes. Wendy Chapman, PhD (Associate Dean of Health Informatics and Chief Learning Health Officer, UT Southwestern Medical Center) asks “Beyond Model Performance: What Kind of Care Does AI Enable?”, arguing that evaluation must extend well beyond accuracy to encompass implementation, human factors, and real-world impact on care. Brian Chapman, Associate Professor at UT Southwestern Medical Center, takes a complementary patient-facing perspective in “‘I’ll Be Your Mirror’: Using NLP to Expand Patient Meaning Making”, challenging the field to build NLP tools that serve the patients who generate the data, not only the institutions that collect it.

\par\vspace{1em}
We are deeply grateful to our program committee, reviewers, and, particularly our shared task organizers for their rigorous and dedicated work. We also thank the annotators who created the shared task datasets, the ACL 2026 organizing team for their support, and every researcher who submitted a paper or participated in a shared task. \#SMM4H-HeaRD would not exist without this community, and we are grateful to see it continue to grow.

\par\vspace{2em}
\#SMM4H-HeaRD 2026 Workshop Chairs and Shared Task Committee
